% Owner: boundary_d. B.4.23 independently passed by root and limits_ab.
\begin{proposition}[B.4.23 (Proved): exceptional-curve nonclosure and persistence in a semistable cap]
Let $\mathcal K^*\to\Delta_s$ be the repaired degree-one
Kato family of B.4.21, with $q=s^{12}$, and let
$\mathcal E\to\Delta_s$ be its unchanged persistent elliptic
axis. For $s\ne0$ let $C_s\simeq\PP^1$ be the exceptional
curve of the single Hopf blow-up.

The union of the $C_s$ has no proper analytic divisor
closure in $\mathcal K^*$. Further finite proper modifications
cannot give these curves a proper analytic divisor closure:
any such closure would push forward to a forbidden one.
In particular, the generic Hopf blowdown does not extend
as a relative proper holomorphic modification, or as a
relative meromorphic contraction represented by a proper
analytic graph, across $s=0$.

There is also a deformation obstruction which does not
assume a particular plumbing. Suppose a smooth complex
fourfold $\mathcal Y$ has holomorphic functions $(Q,r)$,
its divisor $r=0$ is reduced normal crossings, and one
component $K$ contains a neighbourhood of some $C_{s_0}$,
$s_0\ne0$, agreeing with the complementary Kato region
and its ordinary base coordinate $Q$. Suppose this curve
is disjoint from every conductor branch of $r=0$.
Then every sufficiently nearby fibre $(Q,r)=(Q_0,r_0)$
contains an embedded rational curve. Consequently such a
cap cannot have the prescribed ordinary two-tori as its
$r\ne0$ surface fibres. The conclusion applies to the
whole complementary Kato component with these hypotheses;
it makes no assertion about a different central component,
a curve meeting the conductor, or a replacement changing
the holomorphic base map.
\end{proposition}
\begin{proof}
We first use the exact marked blowdown of B.4.10.
The exceptional curve of the single generic Hopf blow-up
meets the persistent elliptic axis transversely at
\[
 [v=-q]\in E_{a(q)}=\C^*/a(q)^{\Z}.
\]
All modifications in B.4.21 are disjoint from that axis.
The intersection section in the repaired family is therefore
still
\[
 p(s)=[-s^{12}].                                      \tag{B.4.23.1}
\]
The units $a(q)$ obey $0<|a(q)|<1$. Choose a holomorphic
logarithm on the full parameter disc and put
\[
 \tau(s)=\frac{\log a(s^{12})}{2\pi i}.
\]
Then $\mathcal E_s=\C/(\Z+\tau(s)\Z)$ and
$\operatorname{Im}\tau(s)>0$. The maps
$(x,y,s)\mapsto[x+y\tau(s)]$, with $(x,y)\in\R^2/\Z^2$,
give a smooth trivialization of the elliptic family.
A logarithmic lift of (B.4.23.1) is
\[
 z(s)=\frac{\log(-s^{12})}{2\pi i}.
\]
On a positive loop around zero, $\tau$ returns to itself
and $z$ changes by $12$. Thus the loop of the section has
class $(12,0)\ne0$ in the integral first homology of this
trivialized family. A continuous section over the disc
would fill that loop. Hence $p$ has no continuous, and
therefore no holomorphic, extension. A meromorphic section
from a nonsingular curve would extend holomorphically:
its proper graph is generically degree one over the curve,
and normalization gives a finite degree-one map to that
normal base. This is equally impossible. A finite base
change of degree $d>0$ multiplies the winding to $(12d,0)$
and does not change the conclusion.

Suppose an analytic divisor $Z$ in $\mathcal K^*$, proper
over the disc, restricted over the punctured disc to the
union of the $C_s$. Its horizontal irreducible part does
not contain $\mathcal E$, since the generic intersection
is the one transverse point (B.4.23.1). Both are Cartier
locally in the smooth threefold, so their intersection
is an analytic curve, proper over the disc. Its reduced
horizontal part has generic degree one. An irreducible
horizontal analytic curve cannot have a positive-dimensional
fibre without being contained in that fibre. Hence its
normalization is finite of degree one over the disc and
is isomorphic to the disc. It extends $p(s)$, a
contradiction. Vertical components of the intersection
have no effect on this horizontal component.

If an additional finite proper modification admitted a
proper analytic closure, its proper image would give
such a $Z$, by the proper mapping theorem; the image has
dimension two and the same generic exceptional curves.
The horizontal exceptional locus of a relative proper
holomorphic blowdown likewise gives $Z$. For a meromorphic
contraction with proper analytic graph, resolve the graph
and apply the same argument to its horizontal exceptional
locus and its proper image in $\mathcal K^*$.
This is a direct marked-intersection proof of a familiar
nonproperness phenomenon for exceptional curves in GSS
surface deformations; see
\citet[Section 1.2]{DlousskyTeleman2012AB} for the primary
context. No classification of all central curves is used.

We next prove the deformation assertion. Write
$C=C_{s_0}$ and $S=K\cap\{Q=Q_0\}$. Near this curve,
$S$ is a smooth surface and $C$ is its exceptional curve,
so
\[
 N_{C/S}\simeq\OO_{\PP^1}(-1).
\]
The base coordinate $Q$ is submersive on $K$ near $C$.
Because $C$ misses every other central branch, the reduced
function $r$ is a local defining function for $K$ near
$C$ and its normal derivative is nowhere zero.
Consequently the differential of $(Q,r)$ gives the exact
normal sequence
\[
 0\longrightarrow\OO_{\PP^1}(-1)
 \longrightarrow N_{C/\mathcal Y}
 \xrightarrow{\ (dQ,dr)\ }
 \OO_{\PP^1}^{\oplus2}\longrightarrow0.
                                                        \tag{B.4.23.2}
\]
Its extension class vanishes because
$H^1(\PP^1,\OO(-1))=0$. Hence
\[
 N_{C/\mathcal Y}\simeq
 \OO(-1)\oplus\OO^{\oplus2},\qquad
 H^1(C,N_{C/\mathcal Y})=0,\qquad
 h^0(C,N_{C/\mathcal Y})=2.                             \tag{B.4.23.3}
\]
Moreover (B.4.23.2) induces an isomorphism on global
sections, since both $H^0(\OO(-1))$ and $H^1(\OO(-1))$
vanish.

Kodaira's embedded-deformation theorem therefore gives a
smooth two-dimensional local parameter space $T$ of
compact embedded deformations of $C$, with tangent space
$H^0(C,N_{C/\mathcal Y})$. In precisely this relative
surface setting one may apply the directly opened
\citet[Proposition 2.3 and its proof, pp.~6--7 of the
original preprint]{DlousskyTeleman2012AB}: after shrinking
the neighbourhood of $C$, $(Q,r)$ is a holomorphic
submersion of relative dimension two, and $C$ is a
compact simply exceptional divisor of its fibre.
That proposition does not require the submersion to be
proper. The general theorem is due to
\citet{Kodaira1962D}; its exact statement is also recorded
in the opening of \citet{Merkulov1995D}.
Shrink this family so its curves
remain in the chosen neighbourhood. Their genus remains
zero, so all are isomorphic to $\PP^1$.
The functions $Q,r$ are constant on every compact
connected curve and thus define a holomorphic map
\[
 T\longrightarrow\C^2.
\]
Its derivative at $[C]$ is precisely the isomorphism on
sections induced by (B.4.23.2). The holomorphic inverse
function theorem shows that the map covers a neighbourhood
of $(Q_0,0)$. In particular the nearby fixed pairs with
$r\ne0$ contain rational curves. A holomorphic map from
$\PP^1$ to a complex torus lifts to its vector-space
universal cover and is constant. Thus those fibres cannot
be the prescribed ordinary two-tori.

The conductor condition is essential. For example, if a
further central component meets $K$ along the persistent
elliptic axis, then its intersection number with $C$ is
one and $N_{K/\mathcal Y}|_C\simeq\OO(-1)$.
The normal sequence need no longer supply the $r$ direction
in (B.4.23.2). The proposition therefore does not exclude
such a changed central arrangement; its proper construction
and comparison with the prescribed family remain separate
requirements.
\end{proof}
